\section{Discussion and Conclusions}
\label{sec:discussion}
\subsection{Main findings}
\label{sec:representation_findings}
Egocentric recordings can be organized into a shared representation of manual work through generated descriptions, shared factory labels, and interval-event extraction. For \textbf{RQ1}, the released logs satisfy the reported field, ordering, and vocabulary checks. For \textbf{RQ2}, 36 of 42 sampled events were fully supported by their generated descriptions, while accuracy against the videos remains unmeasured. These results distinguish structural consistency and textual support from correct observation.

For \textbf{RQ3}, the logs support descriptions of work composition, local sequences, and worker profiles. They quantify behaviors described in the reports and identify questions for further investigation. For \textbf{RQ4}, the examples show how representation choices affect these measurements. In factory \textbf{003}, the process map counts recorded transport under other categories. In the foundry example, an inferred boundary divides a described episode between activities. Synthetic clip joins can add transitions. Interpreting a result therefore requires its category rule, event boundaries, case definition, and denominator.

\subsection{Limits and responsible use}
\label{sec:limits}
The vocabulary is induced from the same factory corpus used for assessment, so the results do not establish transfer to unseen factories. The single-author audit uses generated text and cannot reveal visual errors shared by the generated transcripts and logs.

Missing source durations prevent us from measuring footage coverage or the amount of missing time. Missing selection records prevent us from establishing how representative the processed recordings are. Missing historical execution details prevent exact regeneration of the model outputs. The released artifacts still permit reproduction of the descriptive calculations.

Different task mixes, observation windows, and category assignments can explain differences between worker profiles. Product identifiers and production outcomes are absent, and no intervention was evaluated. The results consequently support hypotheses about process design, not causal diagnoses, efficiency rankings, or estimates of recoverable capacity.

\label{sec:ethics}
Worker identifiers are pseudonymous and may still enable identification when combined with local information~\cite{mannhardt2018privacy}. Use should remain tied to an agreed analytical purpose, with suitable aggregation and review by the people whose work is represented. These inferred profiles do not support employee appraisal, pay, or disciplinary decisions. The source card does not document the original consent procedure; this secondary study collected no new consent.

\subsection{Next steps}
The linked generated transcripts, summaries, vocabularies, and logs provide a basis for testing alternative event representations. Priorities are independent video assessment with multiple evaluators, preservation of source durations, and explicit links from events to supporting generated transcript spans. Product references and real calendar times would be needed to extend the resource to synchronized production-flow analysis.
